% English version of Sec. 3.2 (objective) of reference_kto_formulation.tex
% Build: xelatex reference_kto_objective_en.tex
\documentclass[11pt]{article}
\usepackage[a4paper,margin=2.0cm]{geometry}
\usepackage{amsmath,amssymb,bm}
\usepackage{booktabs}
\usepackage{enumitem}
\usepackage[hidelinks]{hyperref}
\setlength{\parskip}{4pt}
\setlength{\parindent}{0pt}
\DeclareMathOperator{\diag}{diag}

\title{Kinematic Trajectory Optimization of a Reference Clip\\[3pt]\large Objective function}
\author{Companion to \texttt{reference\_kto\_formulation.tex}, Sec.~3.2}
\date{2026-09-04}

\begin{document}
\maketitle

\section*{Objective}

The whole-clip problem minimizes
\begin{equation}
J(\bm Q)=\sum_{t=0}^{N-1} f_t(\bm q_t)\;+\;w_a\sum_{t}\bigl\|(\bm q_{t+1}\ominus\bm q_t)-(\bm q_t\ominus\bm q_{t-1})\bigr\|^2
\label{eq:obj}
\end{equation}
over the trajectory $\bm Q=(\bm q_0,\dots,\bm q_{N-1})$, subject to the hard
constraints (C1)--(C8) of the full formulation. The first term is the
per-frame intent cost, the second the squared discrete acceleration
(smoothness); $\ominus$ is the tangent-space difference (logarithm map on the
base rotation, plain subtraction elsewhere).

Let
\begin{equation}
\mathcal S_t=\bigl\{h:\ \sigma_{h,t}=0,\ (h,t)\notin\mathcal T_{\mathrm{rake}}\bigr\}
\end{equation}
denote the pads that are released and not raking at frame $t$. The per-frame
intent cost is
\begin{equation}
\begin{aligned}
f_t(\bm q)=\;&\sum_{h\in\mathcal S_t}\bigl\|W^p_h\bigl(\bm p_h(\bm q)-\bm p^*_h(t)\bigr)\bigr\|^2
  && \text{swing-pad arc}\\
  &+\sum_{h\in\mathcal S_t} w^r_h\bigl\|\bm e_r\bigl(R_h(\bm q),R^*_h(t)\bigr)\bigr\|^2
  && \text{swing-pad orientation}\\
  &+\sum_{h:\ \beta_{h,t}=1}\bigl\|W^{\mathrm{toe}}\bigl(\bm p_{\mathrm{toe},h}(\bm q)-\bm p^*_{\mathrm{toe},h}(t)\bigr)\bigr\|^2
  && \text{toe brace}\\
  &+\bigl\|W^b\bigl(\bm p_b-\bm p^*_b(t)\bigr)\bigr\|^2+w^b_r\bigl\|\bm e_r\bigl(R_b,R^*_b(t)\bigr)\bigr\|^2
  && \text{base intent}\\
  &+\bigl\|W_\theta\bigl(\bm\theta-\bm\theta^*(t)\bigr)\bigr\|^2
  && \text{posture reference}\\
  &+\bigl\|W_{\mathrm{lim}}(\bm\theta)\bigl(\bm\theta-\bm\theta_{\mathrm{bnd}}\bigr)\bigr\|^2
  && \text{limit-margin spring}\\
  &+w_\gamma^2\|\bm\gamma\|^2 .
  && \text{gimbal return}
\end{aligned}
\label{eq:ft}
\end{equation}
Here $\bm p_h(\bm q),R_h(\bm q)$ are the pad-frame position and orientation
from forward kinematics, $\bm e_r(R,R^*)=\tfrac12\sum_{k=1}^{3}\bm r_k\times\bm r^*_k$
is the world-frame small-angle orientation error, $\bm p_b,R_b$ the base
pose, $\bm\theta$ the joint angles, $\bm\gamma$ the gimbal angles, and
$\beta_{h,t}$ the toe-brace schedule.

\subsection*{Term by term}
\begin{itemize}[leftmargin=1.6em]
\item \textbf{Swing pad.} For every pad that is released and not yet raking,
      the pad follows the Bezier approach arc $\bm p^*_h(t)$ toward its next
      anchor. $W^p_h$ is a $3\times3$ diagonal weight; $w^r_h$ ramps with
      arc progress (weak at the start of a hand swing, strong toward
      landing). This is an intent, not a fact: the arc encodes a preference
      (clear the wall, never drag unloaded spine tips across it), not a
      physical constraint.
\item \textbf{Toe brace.} $W^{\mathrm{toe}}=\diag(300,30,30)$: only the
      wall-normal $x$ component is enforced strongly; $y$ and $z$ ride
      freely, because the 2-DoF ankle cannot place the toe tip exactly in
      3D.
\item \textbf{Base intent.} $\bm p^*_b(t)$ blends linearly from the
      segment-start position to ``press $3$\,cm toward the wall, shift
      $2.5$\,cm toward the supporting side, rise half a step''; $R^*_b(t)$
      interpolates back to upright. The weights are small ($30$--$100$ for
      position, $2$--$50$ for orientation), so the welds decide where the
      base can actually go.
\item \textbf{Posture reference.} $\bm\theta^*(t)$ blends from the
      segment-start angles to the stance angles with the smallest weight
      ($0.3$); it only keeps directions that no other term constrains from
      wandering.
\item \textbf{Limit-margin spring.} $W_{\mathrm{lim}}(\bm\theta)$ is
      nonzero only for joints inside the band within $\mu_s=10\%$ of their
      range from either stop; it pulls such a joint back to the band
      boundary $\bm\theta_{\mathrm{bnd}}$. It differs from the hard margin
      $\bm m$ of constraint (C4): the hard margin is a wall, the spring is
      an early nudge.
\item \textbf{Gimbal return.} $w_\gamma$ makes the solver prefer the arm
      joints over the passive gimbal when satisfying a weld; the gimbal
      turns only when the arm cannot.
\item \textbf{Smoothness.} $w_a$ in \eqref{eq:obj} penalizes the discrete
      acceleration; constraint (C7) bounds it, $w_a$ makes the trajectory
      smoother than the bound alone requires.
\end{itemize}

\subsection*{Default weights}
\begin{center}
\begin{tabular}{@{}lll@{}}
\toprule
Weight & Value & Role \\
\midrule
$W^p_h$ (foot) & $\diag(1000,1000,1000)$ & swing-foot position \\
$w^r_h$ (foot) & $100$ & swing-foot orientation \\
$W^p_h$ (hand) & $500\to2000$ with arc progress & swing-hand position \\
$w^r_h$ (hand) & $50\to500$ with arc progress & swing-hand orientation \\
$W^{\mathrm{toe}}$ & $\diag(300,30,30)$ & toe brace \\
$W^b$, $w^b_r$ & $\diag(30\ldots100)$, $2\ldots50$ & base intent \\
$W_\theta$ & $0.3$ (right-hand swing $10$) & posture reference \\
$W_{\mathrm{lim}}$, $\mu_s$ & $3$, $0.10$ & limit-margin spring \\
$w_\gamma$ & $3$ & gimbal return \\
$w_a$ & to be tuned (order $1$, in rad/s$^2$) & smoothness \\
\bottomrule
\end{tabular}
\end{center}
The swing-hand weights follow
$w=w_{\min}+(w_{\max}-w_{\min})\max\bigl(1-\min(s_a/0.3,1),\,s_a^2\bigr)$
with arc progress $s_a$: strong at the segment start, weak mid-swing, strong
again before landing, which spreads the re-equilibration caused by releasing a
weld over the whole swing instead of a single frame.

\end{document}
